% =========================================================================
\section{Introducción: el problema de calcular el potencial}\label{sec:intro}
% =========================================================================
En el tema anterior~\cite{TEMA1} obtuvimos las leyes que permiten calcular el campo y el potencial creados por
\emph{cualquier} distribución de carga en \emph{cualquier} punto del espacio. Conviene tenerlas a la vista,
porque este tema no añade ninguna ley nueva: se ocupa de cómo \emph{evaluar} en la práctica una de ellas.

% -------------------------------------------------------------------------
\subsection{Las tres descripciones de la electrostática}\label{sec:triangulo}
% -------------------------------------------------------------------------
La densidad de carga $\rho$, el potencial $\phi$ y el campo $\vE$ son tres descripciones equivalentes de un mismo problema
electrostático: conocida una de ellas, las otras dos quedan determinadas (\cref{fig:triangulo}).

\begin{proposicion}[Relaciones fundamentales de la electrostática~\cite{TEMA1}]\label{prop:triangulo}
Sea $\rho(\rp)$ una distribución de carga localizada en un volumen $V'$, y tomemos el origen de potencial en el infinito. Entonces,
para todo punto campo $\vc{r}$,
\begin{align}
  \phi(\vc{r})&=\frac{1}{4\pi\eo}\int_{V'}\frac{\rho(\rp)}{R}\,\dVp, &
  \vE(\vc{r})&=\frac{1}{4\pi\eo}\int_{V'}\frac{\rho(\rp)}{R^{3}}\,\vc{R}\,\dVp,\label{eq:integrales}\\[2pt]
  \vE(\vc{r})&=-\nab\phi(\vc{r}), &
  \phi(\vc{r})-\phi(\vc{r}_0)&=-\int_{\vc{r}_0}^{\vc{r}}\vE\cdot\dl,\label{eq:E-phi}\\[2pt]
  \nab\cdot\vE&=\frac{\rho}{\eo}, &
  \Lap\phi&=-\frac{\rho}{\eo},\label{eq:poisson}
\end{align}
donde $\vc{R}=\vc{r}-\rp$ y $R=\norm{\vc{R}}$.
\end{proposicion}

% =========================================================================
% Figura: el «triángulo» de la electrostática (apuntes manuscritos, pág. 1).
% Vértices rho, phi, E; en cada lado, la relación que permite pasar de una
% descripción a otra (integración directa, derivación o ecuación diferencial).
% Las flechas de ida van por fuera del triángulo y las de vuelta por dentro,
% con los rótulos en el lado correspondiente para que no se superpongan.
% =========================================================================
\begin{figure}[ht]
\centering
\begin{tikzpicture}[x=1cm, y=1cm,
    vert/.style={draw=magacero, line width=0.6pt, fill=magfondo, rounded corners=2pt,
                 inner sep=5pt, minimum width=1.45cm, minimum height=0.75cm, font=\normalsize},
    rel/.style={font=\footnotesize, text=black!85, inner sep=1.5pt},
    ida/.style={-{Stealth[length=5pt,width=4pt]}, line width=0.7pt, magacero, shorten >=17pt, shorten <=17pt},
    vuelta/.style={-{Stealth[length=5pt,width=4pt]}, line width=0.7pt, grispizarra, shorten >=17pt, shorten <=17pt}]
  \coordinate (A) at (0,4.2);   % rho
  \coordinate (B) at (-3.8,0);  % phi
  \coordinate (C) at (3.8,0);   % E
  % normales unitarias exteriores de los lados izquierdo y derecho
  \coordinate (nI) at (-0.742,0.671);
  \coordinate (nD) at (0.742,0.671);
  % lado izquierdo: rho -> phi (por fuera), phi -> rho (por dentro)
  \draw[ida]    ($(A)+0.13*(nI)$) -- ($(B)+0.13*(nI)$);
  \draw[vuelta] ($(B)-0.13*(nI)$) -- ($(A)-0.13*(nI)$);
  \node[rel, anchor=south east] at ($(A)!0.5!(B)+0.25*(nI)$)
       {$\displaystyle\phi=\frac{1}{4\pi\eo}\int_{V'}\frac{\rho}{R}\,\dVp$};
  \node[rel, anchor=north west] at ($(A)!0.57!(B)-0.2*(nI)$) {$\Lap\phi=-\dfrac{\rho}{\eo}$};
  % lado derecho: rho -> E (por fuera), E -> rho (por dentro)
  \draw[ida]    ($(A)+0.13*(nD)$) -- ($(C)+0.13*(nD)$);
  \draw[vuelta] ($(C)-0.13*(nD)$) -- ($(A)-0.13*(nD)$);
  \node[rel, anchor=south west] at ($(A)!0.5!(C)+0.25*(nD)$)
       {$\displaystyle\vE=\frac{1}{4\pi\eo}\int_{V'}\frac{\rho\,\vc{R}}{R^{3}}\,\dVp$};
  \node[rel, anchor=north east] at ($(A)!0.57!(C)-0.2*(nD)$) {$\nab\cdot\vE=\dfrac{\rho}{\eo}$};
  % base: phi -> E (derivación) y E -> phi (integral de línea)
  \draw[ida]    ($(B)+(0,0.13)$) -- ($(C)+(0,0.13)$);
  \draw[vuelta] ($(C)-(0,0.13)$) -- ($(B)-(0,0.13)$);
  \node[rel, anchor=south] at (0,0.22) {$\vE=-\nab\phi$};
  \node[rel, anchor=north] at (0,-0.22) {$\displaystyle\phi(\vc{r})-\phi(\vc{r}_0)=-\int_{\vc{r}_0}^{\vc{r}}\vE\cdot\dl$};
  % vértices
  \node[vert] at (A) {$\rho(\rp)$};
  \node[vert] at (B) {$\phi(\vc{r})$};
  \node[vert] at (C) {$\vE(\vc{r})$};
\end{tikzpicture}
\caption{Las tres descripciones equivalentes de un problema electrostático (\cref{prop:triangulo}). Las flechas magenta (exteriores) se recorren
\emph{integrando} sobre las fuentes o \emph{derivando}; las grises (interiores), resolviendo una ecuación diferencial o integrando a lo largo de un
camino. El desarrollo multipolar aproxima la flecha $\rho\to\phi$, que es la más difícil de recorrer en la práctica.}
\label{fig:triangulo}
\end{figure}


\begin{observacion}[Por qué preferimos el potencial]\label{obs:potencial-escalar}
De las dos integrales de \cref{eq:integrales}, la del potencial es siempre la más sencilla: su integrando es escalar (un número por
punto en lugar de tres) y decae como $1/R$ en lugar de como $1/R^{2}$. Por eso, en todo el tema calcularemos primero el potencial y
obtendremos después el campo derivando, $\vE=-\nab\phi$.
\end{observacion}

% -------------------------------------------------------------------------
\subsection{Planteamiento: una aproximación a grandes distancias}\label{sec:planteamiento}
% -------------------------------------------------------------------------
Aun así, la integral
\begin{equation}\label{eq:pot-integral}
  \phi(\vc{r})=\frac{1}{4\pi\eo}\int_{V'}\frac{\rho(\rp)}{\norm{\vc{r}-\rp}}\,\dVp
\end{equation}
es, en general, \emph{difícil de calcular en la práctica}: solo para distribuciones muy simétricas (esferas, cilindros, planos) admite una
expresión cerrada, y en los casos más generales ni siquiera conocemos con detalle $\rho$. Lo que sí sabemos es que si nos alejamos
suficientemente de la distribución, sus detalles dejan de importar y toda ella se ve como una \emph{carga puntual}. Entre ambos extremos
(conocer $\rho$ punto a punto o verla como un punto) hay un camino intermedio: caracterizar la distribución mediante unas pocas
cantidades, que dependen solo de \emph{algunos} detalles de la misma, y que bastan para obtener el potencial con la precisión deseada.

\begin{definicion}[Momentos multipolares]\label{def:multipolos}
Llamamos \emph{momentos multipolares} (o \emph{multipolos eléctricos}) de una distribución de carga a las cantidades
\[
  \int_{V'}\rho(\rp)\,\dVp,\qquad \int_{V'}\rho(\rp)\,\rp\,\dVp,\qquad \int_{V'}\rho(\rp)\bigl(3x'_jx'_k-\delt{j}{k}r'^{2}\bigr)\dVp,\ \dots
\]
(carga neta, momento dipolar, tensor cuadrupolar, \dots), que solo dependen de la distribución y del origen elegido, pero \emph{no} del
punto campo. El \emph{desarrollo multipolar} es la expresión del potencial como una serie cuyos términos son productos de un momento
multipolar por una función sencilla del punto campo, y que decaen cada vez más deprisa con la distancia.
\end{definicion}

El desarrollo permite así sustituir una integral distinta para cada punto campo por unas pocas integrales, calculadas una sola vez, y
\emph{sacar} las coordenadas del punto campo fuera de la integral. Construiremos el desarrollo en dos etapas. Primero estudiaremos el sistema más
sencillo de carga neta nula, el \emph{dipolo eléctrico} (\cref{sec:dipolo}), y su generalización a distribuciones continuas de dipolos, la
\emph{polarización} (\cref{sec:polarizacion}). Después desarrollaremos $1/R$ en serie de potencias de $r'/r$ y obtendremos el desarrollo
completo (\cref{sec:multipolar}), cuyo error de truncamiento acotaremos en la \cref{sec:error}. Los seminarios (\cref{sec:seminarios}) aplican
todo ello a los problemas de la lista del tema~\cite{PROBLEMAS-T2}. El desarrollo sigue los apuntes de teoría de la asignatura~\cite{APUNTES-T2}; como
bibliografía de apoyo, el dipolo eléctrico se trata en~\cite[cap.~3]{GRIFFITHS}, \cite[cap.~8]{WANGSNESS} y~\cite[cap.~4]{VICTORIANO}, y las expresiones multipolares en las
mismas referencias y, con más profundidad, en~\cite{JACKSON}.

% -------------------------------------------------------------------------
\subsection{Notación y convenios del tema}\label{sec:notacion}
% -------------------------------------------------------------------------
\begin{notacion}[Puntos campo y fuente]\label{not:puntos}
Seguimos la notación de los temas anteriores~\cite{TEMA0,TEMA1}, que a su vez adopta la de~\cite{NAVAS}:
\begin{itemize}
  \item $\vc{r}=(x,y,z)$ es el \emph{punto campo}, donde se evalúan $\phi$ y $\vE$, con coordenadas esféricas $(r,\theta,\varphi)$ y
        versores $\{\ur,\uth,\uph\}$; las magnitudes primadas, $\rp=(x',y',z')$, $r'=\norm{\rp}$, $\dVp$, $\nabp$, se refieren
        al \emph{punto fuente}.
  \item $\vc{R}=\vc{r}-\rp$ es el vector que va de la fuente al punto campo, $R=\norm{\vc{R}}$ y $\uR=\vc{R}/R$.
  \item $\theta'$ es el ángulo que forman $\vc{r}$ y $\rp$, de modo que $\vc{r}\cdot\rp=rr'\cosp$.
  \item $\cd{x},\cd{y},\cd{z}$ son los \emph{cosenos directores} del punto campo, es decir, las componentes cartesianas de $\ur$:
        \[
          \ur=\cd{x}\,\ux+\cd{y}\,\uy+\cd{z}\,\uz,\qquad
          \cd{x}=\frac{x}{r}=\sen\theta\cos\varphi,\quad \cd{y}=\frac{y}{r}=\sen\theta\sen\varphi,\quad \cd{z}=\frac{z}{r}=\cos\theta,
        \]
        que verifican $\cd{x}^{2}+\cd{y}^{2}+\cd{z}^{2}=1$.
  \item Los índices latinos $j,k\in\{x,y,z\}$ recorren las componentes cartesianas; $x'_j$ denota $x',y',z'$, y $\delt{j}{k}$ es la delta de
        Kronecker. Las sumas se escriben siempre de forma explícita, $\sum_{j,k}$.
\end{itemize}
\end{notacion}

\begin{observacion}[Tres usos de la letra $\ell$]\label{obs:ell}
Por fidelidad al formulario del curso~\cite{FORMULARIO}, la letra $\ell$ aparece en tres papeles que el contexto distingue sin ambigüedad:
como \emph{índice} del desarrollo ($\ell=0,1,2,\dots$, orden del multipolo), como \emph{vector} $\vl$ de separación de las cargas de un
dipolo (siempre con flecha) y como \emph{coseno director} $\cd{j}$ (siempre con subíndice cartesiano). En los seminarios, además, $\ell$ es
la longitud de una barra.
\end{observacion}

Como en el Tema~1, la constante de Coulomb es $1/(4\pi\eo)\simeq8{,}988\cdot10^{9}\,\mathrm{N\,m^{2}\,C^{-2}}$~\cite{CODATA} y todas las
distribuciones se suponen \emph{localizadas} (contenidas en una esfera de radio finito), de modo que el potencial se anula en el infinito.
